% Compile: pdflatex report.tex  (chạy 2 lần)
\documentclass[12pt, a4paper]{article}

% ── Encoding & Tiếng Việt ────────────────────────────────────────
\usepackage[utf8]{inputenc}
\usepackage[T5]{fontenc}
\usepackage[vietnamese]{babel}

% ── Bố cục trang ─────────────────────────────────────────────────
\usepackage[top=2.5cm, bottom=2.5cm, left=3cm, right=2.5cm]{geometry}
\usepackage{setspace}
\onehalfspacing

% ── Toán học ─────────────────────────────────────────────────────
\usepackage{amsmath, amssymb}

% ── Bảng ─────────────────────────────────────────────────────────
\usepackage{booktabs}
\usepackage{multirow}
\usepackage{array}
\usepackage{longtable}
\newcolumntype{C}[1]{>{\centering\arraybackslash}p{#1}}

% ── Hình ảnh ─────────────────────────────────────────────────────
\usepackage{graphicx}
\usepackage{xcolor}

% ── Khác ─────────────────────────────────────────────────────────
\usepackage{hyperref}
\usepackage{caption}
\usepackage{subcaption}
\usepackage{enumitem}
\usepackage{float}

\hypersetup{colorlinks=true, linkcolor=blue, citecolor=blue, urlcolor=blue}

% ─────────────────────────────────────────────────────────────────
\begin{document}

% ── Trang bìa ────────────────────────────────────────────────────
\begin{titlepage}
  \centering
  % Đặt file logo (vd: logo.png) cùng thư mục với report.tex rồi bỏ comment dòng dưới
  % \includegraphics[width=3.5cm]{logo.png}\\[0.8cm]
  {\Large\bfseries TÊN TRƯỜNG ĐẠI HỌC}\\[0.3cm]
  {\large Khoa Công nghệ Thông tin}\\[3cm]

  \rule{\linewidth}{0.8pt}\\[0.5cm]
  {\LARGE\bfseries Ứng dụng Học Tăng cường vào Thuật toán\\[0.3cm]
   Tìm kiếm Lân cận Diện rộng Thích ứng\\[0.3cm]
   cho Bài toán Thiết kế Lãnh thổ}\\[0.4cm]
  {\large\itshape Báo cáo Thực nghiệm}\\[0.3cm]
  \rule{\linewidth}{0.8pt}\\[2.5cm]

  \begin{minipage}{0.48\textwidth}
    \begin{flushleft}
      \textbf{Tác giả:}\\[0.2cm]
      Trần Đức Thái\\
      \texttt{thaitd@vietnamlab.vn}
    \end{flushleft}
  \end{minipage}

  \vfill
  {\large Tháng 7, 2026}
\end{titlepage}

% ── Tóm tắt (trang riêng) ────────────────────────────────────────
\newpage
\thispagestyle{plain}
\begin{abstract}
Báo cáo này trình bày kết quả thực nghiệm của thuật toán \textbf{ALNS kết hợp
Học Tăng cường dạng Repair (ALNS+RL Repair)} áp dụng cho Bài toán Thiết kế Lãnh thổ
(District and Territory Design Problem -- DTDP).
Tác nhân RL sử dụng mô hình Q-learning tuyến tính để học chính sách gán node vào
territory trong giai đoạn \emph{repair} của Adaptive Large Neighbourhood Search (ALNS).
Thực nghiệm được thực hiện trên 120 đồ thị thuộc hai họ: TGraph (đồ thị phẳng ngẫu
nhiên, 500--700 nút) và GGraph (đồ thị lưới $27\!\times\!27$ đến $33\!\times\!33$,
ba cấu hình depot).
Kết quả cho thấy ALNS+RL Repair tìm được nghiệm \emph{khả thi hoàn toàn} trên GGraph
($\text{Infeasibility} \approx 0$) nhưng gặp khó khăn về tính khả thi trên TGraph
(infeasibility trung bình $\approx 0{,}21$).
Chất lượng mục tiêu thua VNS baseline khoảng 18\% trên TGraph và 16\% trên GGraph.
Tác nhân generalize tốt về \emph{tính khả thi} trên GGraph nhưng chưa học được chiến
lược tối ưu hóa diameter cạnh tranh với VNS.
\end{abstract}

% ── Mục lục (trang riêng) ────────────────────────────────────────
\newpage
\tableofcontents
\newpage

% =================================================================
\section{Giới thiệu}
% =================================================================

Bài toán Thiết kế Lãnh thổ (DTDP) xuất hiện trong nhiều bài toán thực tiễn như
phân vùng logistics, quản lý hành chính, và thiết kế mạng lưới dịch vụ.
Mục tiêu là phân hoạch tập nút của một đồ thị thành $k$ lãnh thổ cân bằng, sao
cho đường kính tối đa trong các lãnh thổ được tối thiểu hóa.

Các thuật toán metaheuristic như Variable Neighbourhood Search (VNS) và Adaptive
Large Neighbourhood Search (ALNS) đã cho thấy hiệu quả tốt trên bài toán này.
Báo cáo nghiên cứu liệu \textbf{Học Tăng cường (RL)} có thể cải thiện ALNS bằng
cách học chính sách sửa chữa nghiệm (repair) thông minh hơn chiến lược greedy
truyền thống hay không.

Tác nhân RL dùng \emph{Q-learning tuyến tính} để quyết định gán mỗi node bị xóa
vào territory nào trong quá trình repair, thay vì dùng heuristic greedy cố định.
Hai họ benchmark được sử dụng:
\begin{itemize}
  \item \textbf{TGraph} -- 30 đồ thị phẳng ngẫu nhiên (500, 600, 700 nút).
  \item \textbf{GGraph} -- 90 đồ thị lưới (486, 600, 726 nút) với ba cấu hình
        vị trí depot: Center, Corners, Diagonal.
\end{itemize}

% =================================================================
\section{Phát biểu Bài toán}
% =================================================================

\subsection{Định nghĩa DTDP}

Cho đồ thị $G = (V, E)$ với $|V| = n$ nút.
Mỗi nút $v \in V$ mang các thuộc tính: khối lượng công việc $w_v$, số khách
hàng $c_v$, và nhu cầu $d_v$.
DTDP tìm phân hoạch $\Pi = \{D_1, D_2, \ldots, D_k\}$ của $V$ thành $k$ lãnh thổ sao cho:

\begin{enumerate}
  \item \textbf{Tối thiểu hóa} $F(\Pi) = \max_{j}\,\operatorname{diam}(D_j)$,
        đường kính lớn nhất trong các lãnh thổ.
  \item \textbf{Ràng buộc cân bằng}: với mỗi lãnh thổ $D_j$,
        \[
          \left|\frac{w(D_j)}{\bar{w}} - 1\right| \leq \delta,
          \quad \bar{w} = \frac{\sum_v w_v}{k},
        \]
        trong đó $\delta = 0{,}05$ là độ lệch cho phép (5\%).
\end{enumerate}

\subsection{Hàm Merit}

Để cân bằng giữa tối ưu hóa diameter và đảm bảo tính khả thi, hàm merit được
định nghĩa:
\[
  M(\Pi, \lambda) = \lambda \cdot \frac{F(\Pi)}{D_G} + (1-\lambda) \cdot G(\Pi),
\]
trong đó $D_G$ là đường kính của toàn đồ thị, $G(\Pi)$ là tổng vi phạm cân
bằng tải của tất cả các lãnh thổ, và $\lambda = 0{,}7$ là trọng số.

% =================================================================
\section{Phương pháp}
% =================================================================

\subsection{Tổng quan ALNS}

Thuật toán ALNS hoạt động theo vòng lặp \textbf{destroy--repair}:
\begin{enumerate}
  \item \textbf{Destroy}: xóa $k$ node khỏi nghiệm hiện tại theo một trong ba
        chiến lược:
        \begin{itemize}
          \item \emph{Random}: xóa ngẫu nhiên $k$ node.
          \item \emph{Worst Diameter}: xóa các node thuộc đường kính lớn nhất
                và lân cận của chúng.
          \item \emph{Infeasible}: xóa ngẫu nhiên $k$ node từ territory vi phạm
                cân bằng nặng nhất.
        \end{itemize}
  \item \textbf{Repair}: gán lại các node bị xóa vào các territory.
  \item \textbf{Chấp nhận}: dùng tiêu chí Simulated Annealing để chấp nhận hoặc
        từ chối nghiệm mới.
\end{enumerate}

\subsection{RL trong giai đoạn Repair}

Thay vì dùng heuristic greedy cố định, tác nhân RL học cách gán node vào territory
tối ưu. Với mỗi node $v$ cần gán, tác nhân quan sát \textbf{vector đặc trưng} mô
tả trạng thái của từng territory và chọn territory tốt nhất.

\subsubsection{Không gian trạng thái}

Vector trạng thái $s_{v,d}$ mô tả việc gán node $v$ vào territory $d$, gồm 8 đặc trưng:
\[
  s_{v,d} = \left[
    \underbrace{\text{excess}_{d,w},\; \text{excess}_{d,c},\; \text{excess}_{d,q}}_{\text{dư thừa hoạt động (3)}},\;
    \underbrace{\text{diam}_d / D_G}_{\text{đường kính (1)}},\;
    \underbrace{w_v/\bar{w},\; c_v/\bar{c},\; q_v/\bar{q}}_{\text{thuộc tính node (3)}},\;
    \underbrace{\bar{d}(v, D_d)}_{\text{khoảng cách (1)}}
  \right],
\]
trong đó $\text{excess}_{d,a} = (A_{d,a} - \bar{a})/\bar{a}$ là độ dư thừa hoạt động
$a$ của territory $d$.

\subsubsection{Hành động và phần thưởng}

\textbf{Hành động}: chọn index territory $a \in \{0, 1, \ldots, k-1\}$ để gán node.

\textbf{Phần thưởng}: cải thiện tức thời của hàm merit sau khi gán:
\[
  r = M_{\text{trước}} - M_{\text{sau}}.
\]

\subsubsection{Mô hình Q-learning tuyến tính}

Q-value được xấp xỉ tuyến tính: $Q(s_{v,a}, a) = \mathbf{w}_a^\top s_{v,a}$,
trong đó mỗi territory $a$ có vector trọng số $\mathbf{w}_a \in \mathbb{R}^8$ riêng.

Cập nhật TD(0):
\[
  \mathbf{w}_a \leftarrow \mathbf{w}_a + \alpha \left[
    r + \gamma \max_{b} Q(s'_{v',b}, b) - Q(s_{v,a}, a)
  \right] s_{v,a},
\]
với learning rate $\alpha = 0{,}01$, discount factor $\gamma = 0{,}9$,
$\varepsilon$-greedy exploration ($\varepsilon$ giảm từ 1.0 xuống 0.05).

\subsection{Tham số thực nghiệm}

\begin{table}[H]
  \centering
  \caption{Tham số ALNS+RL Repair}
  \begin{tabular}{lll}
    \toprule
    Tham số & Giá trị & Ghi chú \\
    \midrule
    $\lambda$ & 0{,}7 & Trọng số merit \\
    $\delta$ & 0{,}05 & Độ lệch cân bằng cho phép \\
    $k$ (districts) & 10 & Số lãnh thổ \\
    $k_{\text{remove}}$ & 10 & Số node xóa mỗi bước destroy \\
    Số vòng lặp ALNS & 300 & Mỗi instance \\
    $\alpha$ & 0{,}01 & Learning rate \\
    $\gamma$ & 0{,}90 & Discount factor \\
    $\varepsilon_{\min}$ & 0{,}05 & Epsilon tối thiểu \\
    \bottomrule
  \end{tabular}
\end{table}

% =================================================================
\section{Thiết kế Thực nghiệm}
% =================================================================

\subsection{Tập dữ liệu}

\begin{table}[H]
  \centering
  \caption{Tổng quan benchmark}
  \begin{tabular}{llccc}
    \toprule
    Họ & Tập & Kích thước & Instances & Mục đích \\
    \midrule
    \multirow{4}{*}{TGraph} & planar500 & 500 nút & 10 & Train + Eval \\
                             & planar600 & 600 nút & 10 & Eval (size $+$20\%) \\
                             & planar700 & 700 nút & 10 & Eval (size $+$40\%) \\
    \midrule
    \multirow{3}{*}{GGraph 27$\times$27} & Center486   & 486 nút & 10 & Train + Eval \\
                                          & Corners486  & 486 nút & 10 & Eval (depot mới) \\
                                          & Diagonal486 & 486 nút & 10 & Eval (depot mới) \\
    \multirow{3}{*}{GGraph 30$\times$30} & Center600   & 600 nút & 10 & Eval \\
                                          & Corners600  & 600 nút & 10 & Eval \\
                                          & Diagonal600 & 600 nút & 10 & Eval \\
    \multirow{3}{*}{GGraph 33$\times$33} & Center726   & 726 nút & 10 & Eval \\
                                          & Corners726  & 726 nút & 10 & Eval \\
                                          & Diagonal726 & 726 nút & 10 & Eval \\
    \bottomrule
  \end{tabular}
\end{table}

\subsection{Quy trình}

\begin{enumerate}
  \item \textbf{Training TGraph}: train agent tuần tự trên planar500 G0--G9,
        tích lũy kiến thức vào model \texttt{model\_tgraph.npz}.
  \item \textbf{Eval TGraph}: dùng agent đã train, chạy với $\varepsilon =
        \varepsilon_{\min}$ (greedy thuần) trên 30 instances.
  \item \textbf{Training GGraph}: train agent trên Center486 G0--G9, lưu vào
        \texttt{model\_ggraph.npz}.
  \item \textbf{Eval GGraph}: đánh giá 90 instances (3 kích thước $\times$
        3 depot $\times$ 10 instances).
\end{enumerate}

% =================================================================
\section{Kết quả}
% =================================================================

\subsection{TGraph -- Đồ thị phẳng ngẫu nhiên}

\begin{table}[H]
  \centering
  \caption{Kết quả tổng hợp TGraph (trung bình 10 instances)}
  \begin{tabular}{lcccccc}
    \toprule
    \multirow{2}{*}{Nhóm} &
    \multicolumn{2}{c}{ALNS+RL Repair} &
    \multicolumn{2}{c}{VNS Baseline} &
    \multirow{2}{*}{Gap Obj} &
    \multirow{2}{*}{Nhận xét} \\
    \cmidrule(lr){2-3}\cmidrule(lr){4-5}
    & Obj & Inf & Obj & Inf & & \\
    \midrule
    planar500 & 54{,}296 & 0{,}2351 & 47{,}641 & 0{,}0000 & $+$14{,}0\% & In-dist \\
    planar600 & 54{,}729 & 0{,}1967 & 47{,}709 & 0{,}0000 & $+$14{,}7\% & Size $+$20\% \\
    planar700 & 57{,}633 & 0{,}1890 & 46{,}016 & 0{,}0000 & $+$25{,}2\% & Size $+$40\% \\
    \midrule
    \textbf{Tổng} & \textbf{55{,}553} & \textbf{0{,}2069} & \textbf{47{,}122} &
    \textbf{0{,}0000} & \textbf{$+$17{,}9\%} & \\
    \bottomrule
  \end{tabular}
\end{table}

\textbf{Nhận xét:}
\begin{itemize}
  \item ALNS+RL Repair có infeasibility trung bình $\approx 0{,}21$, trong khi VNS
        luôn tìm được nghiệm \emph{khả thi hoàn toàn} ($G = 0$).
        Đây là điểm yếu nghiêm trọng của phương pháp đề xuất trên TGraph.
  \item Objective gap tăng theo kích thước đồ thị: $+14\%$ (500 nút) $\to$
        $+25{,}2\%$ (700 nút), cho thấy agent \emph{không generalize tốt} khi
        graph lớn hơn trên TGraph.
  \item Nguyên nhân có thể do cơ chế repair dạng node-by-node không đảm bảo
        tính liên thông và cân bằng tải đồng thời khi graph có cấu trúc phức tạp.
\end{itemize}

\subsection{GGraph -- Đồ thị lưới}

\begin{table}[H]
  \centering
  \caption{Kết quả tổng hợp GGraph (trung bình 10 instances mỗi nhóm)}
  \begin{tabular}{llccccl}
    \toprule
    Kích thước & Depot &
    \multicolumn{2}{c}{ALNS+RL} &
    \multicolumn{2}{c}{VNS} &
    Gap \\
    \cmidrule(lr){3-4}\cmidrule(lr){5-6}
    & & Obj & Inf & Obj & Inf & Obj \\
    \midrule
    \multirow{3}{*}{27$\times$27 (486)} &
    Center   & 233{,}0 & 0{,}0000 & 218{,}8 & 0{,}0000 & $+$6{,}5\% \\
    & Corners  & 254{,}2 & 0{,}0010 & 201{,}7 & 0{,}0074 & $+$26{,}0\% \\
    & Diagonal & 247{,}0 & 0{,}0000 & 199{,}6 & 0{,}0319 & $+$23{,}7\% \\
    \midrule
    \multirow{3}{*}{30$\times$30 (600)} &
    Center   & 260{,}9 & 0{,}0000 & 238{,}3 & 0{,}0001 & $+$9{,}5\% \\
    & Corners  & 260{,}1 & 0{,}0000 & 216{,}9 & 0{,}0016 & $+$19{,}9\% \\
    & Diagonal & 274{,}8 & 0{,}0039 & 241{,}8 & 0{,}0082 & $+$13{,}6\% \\
    \midrule
    \multirow{3}{*}{33$\times$33 (726)} &
    Center   & 306{,}2 & 0{,}0004 & 240{,}5 & 0{,}0008 & $+$27{,}3\% \\
    & Corners  & 285{,}0 & 0{,}0010 & 249{,}9 & 0{,}0032 & $+$14{,}0\% \\
    & Diagonal & 299{,}7 & 0{,}0000 & 282{,}0 & 0{,}0199 & $+$6{,}3\% \\
    \midrule
    \multicolumn{2}{l}{\textbf{Tổng (90 instances)}} &
    \textbf{269{,}0} & \textbf{0{,}0007} & \textbf{232{,}2} & \textbf{0{,}0081} &
    \textbf{$+$15{,}8\%} \\
    \bottomrule
  \end{tabular}
\end{table}

\textbf{Nhận xét:}
\begin{itemize}
  \item \textbf{Tính khả thi}: ALNS+RL Repair gần như luôn tìm được nghiệm khả thi
        trên GGraph (infeasibility $\approx 0$), thậm chí tốt hơn VNS ở một số nhóm
        (Corners486, Diagonal486, Diagonal600, Diagonal726 có infeasibility VNS cao hơn).
  \item \textbf{Depot trung tâm (Center)}: gap thấp nhất, $+6{,}5\%$ (486 nút) và
        $+9{,}5\%$ (600 nút) -- cho thấy agent generalize tốt về kích thước
        khi cấu hình depot giống tập train.
  \item \textbf{Depot không quen (Corners, Diagonal)}: gap cao hơn ($+14$--$+26\%$),
        đặc biệt Corners486 ($+26\%$) và Center726 ($+27{,}3\%$).
        Agent chưa học được cách xử lý depot ở vị trí cực trị hiệu quả.
  \item Một số instances riêng lẻ ALNS+RL \emph{tốt hơn} VNS (gap âm),
        ví dụ Center486\_G0 ($-6$), Center486\_G2 ($-13$), Diagonal726\_G3 ($-48$).
\end{itemize}

% =================================================================
\section{Thảo luận}
% =================================================================

\subsection{So sánh với VNS}

\begin{table}[H]
  \centering
  \caption{Tổng hợp so sánh ALNS+RL Repair vs VNS}
  \begin{tabular}{lccc}
    \toprule
    Tiêu chí & ALNS+RL Repair & VNS & Nhận xét \\
    \midrule
    Tính khả thi -- TGraph & \textcolor{red}{Kém} (Inf $\approx 0{,}21$) &
    \textcolor{blue}{Tốt} (Inf $= 0$) & VNS tốt hơn \\
    Tính khả thi -- GGraph & \textcolor{blue}{Tốt} (Inf $\approx 0$) &
    \textcolor{blue}{Tốt} & Tương đương \\
    Objective -- TGraph & $+17{,}9\%$ tệ hơn & Baseline & VNS tốt hơn \\
    Objective -- GGraph Center & $+6{,}5$--$9{,}5\%$ tệ hơn & Baseline & Cạnh tranh \\
    Objective -- GGraph Cross & $+13{,}6$--$27{,}3\%$ tệ hơn & Baseline & VNS tốt hơn \\
    \bottomrule
  \end{tabular}
\end{table}

\subsection{Phân tích nguyên nhân}

\paragraph{TGraph -- infeasibility cao.}
Cơ chế repair node-by-node không có cơ chế đảm bảo tính cân bằng tải toàn cục.
Mỗi lần gán, agent tối ưu hóa merit cục bộ nhưng không nhìn xa được hậu quả
của toàn bộ quá trình repair. Trên TGraph với cấu trúc không đều, điều này dẫn
đến mất cân bằng tích lũy.

\paragraph{GGraph -- tốt hơn TGraph.}
Cấu trúc lưới đều đặn của GGraph giúp agent dễ học pattern gán node hơn.
Node trong lưới có vị trí hình học rõ ràng, khiến các đặc trưng khoảng cách
và excess trong state vector trở nên có ý nghĩa hơn.

\paragraph{Generalization về kích thước.}
Trên GGraph Center, gap tăng vừa phải từ 486 ($+6{,}5\%$) lên 600 ($+9{,}5\%$)
nhưng tăng mạnh ở 726 ($+27{,}3\%$). Điều này gợi ý ngưỡng kích thước mà
agent bắt đầu gặp khó khăn khi domain shift quá lớn.

\subsection{Hướng cải tiến}

\begin{enumerate}
  \item \textbf{Thêm local search sau repair}: sau khi RL gán xong, áp dụng
        local search (node best improvement) để sửa vi phạm còn lại.
  \item \textbf{Reward shaping}: thêm penalty nặng cho việc tạo ra infeasibility,
        không chỉ tối thiểu hóa merit.
  \item \textbf{Deeper RL}: thay Q-learning tuyến tính bằng DQN hoặc Actor-Critic
        để xử lý state phức tạp hơn.
  \item \textbf{Train đa dạng}: train trên cả TGraph và GGraph đa dạng hơn
        thay vì chỉ một loại instance mỗi lần.
\end{enumerate}

% =================================================================
\section{Kết luận}
% =================================================================

Báo cáo đã trình bày và đánh giá thuật toán ALNS+RL Repair trên bài toán DTDP.
Kết quả cho thấy:
\begin{itemize}
  \item Phương pháp \textbf{khả thi tốt trên GGraph} nhưng \textbf{gặp khó khăn về
        infeasibility trên TGraph}.
  \item Chất lượng mục tiêu tổng thể còn thua VNS baseline ($+17{,}9\%$ TGraph,
        $+15{,}8\%$ GGraph).
  \item Agent \textbf{generalize tốt về kích thước} khi depot không thay đổi
        (TGraph, GGraph Center), nhưng \textbf{kém hơn khi loại depot thay đổi}
        (Corners, Diagonal).
  \item Tiềm năng: một số instances riêng lẻ ALNS+RL vượt trội VNS, cho thấy
        hướng RL vào repair có triển vọng nếu được cải tiến thêm.
\end{itemize}

% =================================================================
\appendix
\section{Kết quả chi tiết TGraph}
% =================================================================

\begin{longtable}{lccccc}
  \caption{Kết quả từng instance TGraph} \\
  \toprule
  Instance & ALNS Obj & ALNS Inf & VNS Obj & VNS Inf & $\Delta$ Obj \\
  \midrule
  \endfirsthead
  \toprule
  Instance & ALNS Obj & ALNS Inf & VNS Obj & VNS Inf & $\Delta$ Obj \\
  \midrule
  \endhead
  \bottomrule
  \endfoot
  planar500\_G0 & 52{,}592 & 0{,}1779 & 46{,}111 & 0{,}0000 & $+$6{,}48 \\
  planar500\_G1 & 52{,}378 & 0{,}1786 & 45{,}341 & 0{,}0000 & $+$7{,}04 \\
  planar500\_G2 & 64{,}797 & 0{,}1835 & 49{,}825 & 0{,}0000 & $+$14{,}97 \\
  planar500\_G3 & 64{,}218 & 0{,}4112 & 48{,}767 & 0{,}0000 & $+$15{,}45 \\
  planar500\_G4 & 52{,}138 & 0{,}2180 & 45{,}008 & 0{,}0000 & $+$7{,}13 \\
  planar500\_G5 & 50{,}198 & 0{,}1564 & 47{,}548 & 0{,}0000 & $+$2{,}65 \\
  planar500\_G6 & 50{,}251 & 0{,}2364 & 47{,}225 & 0{,}0000 & $+$3{,}03 \\
  planar500\_G7 & 50{,}970 & 0{,}3479 & 48{,}005 & 0{,}0000 & $+$2{,}96 \\
  planar500\_G8 & 53{,}337 & 0{,}1623 & 49{,}316 & 0{,}0000 & $+$4{,}02 \\
  planar500\_G9 & 52{,}080 & 0{,}2785 & 49{,}263 & 0{,}0000 & $+$2{,}82 \\
  \midrule
  planar600\_G0 & 51{,}886 & 0{,}0509 & 45{,}523 & 0{,}0000 & $+$6{,}36 \\
  planar600\_G1 & 53{,}974 & 0{,}1068 & 50{,}656 & 0{,}0000 & $+$3{,}32 \\
  planar600\_G2 & 56{,}591 & 0{,}1498 & 44{,}752 & 0{,}0000 & $+$11{,}84 \\
  planar600\_G3 & 53{,}876 & 0{,}2718 & 47{,}073 & 0{,}0000 & $+$6{,}80 \\
  planar600\_G4 & 52{,}564 & 0{,}2855 & 47{,}281 & 0{,}0000 & $+$5{,}28 \\
  planar600\_G5 & 56{,}406 & 0{,}2920 & 48{,}253 & 0{,}0000 & $+$8{,}15 \\
  planar600\_G6 & 54{,}664 & 0{,}1802 & 49{,}483 & 0{,}0000 & $+$5{,}18 \\
  planar600\_G7 & 60{,}130 & 0{,}1217 & 49{,}206 & 0{,}0000 & $+$10{,}92 \\
  planar600\_G8 & 55{,}234 & 0{,}1793 & 46{,}365 & 0{,}0000 & $+$8{,}87 \\
  planar600\_G9 & 51{,}964 & 0{,}3286 & 48{,}493 & 0{,}0000 & $+$3{,}47 \\
  \midrule
  planar700\_G0 & 58{,}625 & 0{,}2217 & 46{,}472 & 0{,}0000 & $+$12{,}15 \\
  planar700\_G1 & 58{,}179 & 0{,}1123 & 44{,}097 & 0{,}0000 & $+$14{,}08 \\
  planar700\_G2 & 50{,}747 & 0{,}0795 & 44{,}817 & 0{,}0000 & $+$5{,}93 \\
  planar700\_G3 & 55{,}636 & 0{,}1540 & 46{,}603 & 0{,}0000 & $+$9{,}03 \\
  planar700\_G4 & 67{,}832 & 0{,}3059 & 46{,}419 & 0{,}0000 & $+$21{,}41 \\
  planar700\_G5 & 63{,}709 & 0{,}3066 & 45{,}209 & 0{,}0000 & $+$18{,}50 \\
  planar700\_G6 & 57{,}558 & 0{,}2162 & 49{,}892 & 0{,}0000 & $+$7{,}67 \\
  planar700\_G7 & 52{,}436 & 0{,}2234 & 44{,}879 & 0{,}0000 & $+$7{,}56 \\
  planar700\_G8 & 55{,}607 & 0{,}1521 & 45{,}113 & 0{,}0000 & $+$10{,}49 \\
  planar700\_G9 & 56{,}006 & 0{,}1185 & 46{,}656 & 0{,}0000 & $+$9{,}35 \\
\end{longtable}

% =================================================================
\section{Kết quả chi tiết GGraph}
% =================================================================

\begin{longtable}{lccccc}
  \caption{Kết quả từng instance GGraph 27$\times$27 (486 nút)} \\
  \toprule
  Instance & ALNS Obj & ALNS Inf & VNS Obj & VNS Inf & $\Delta$ Obj \\
  \midrule
  \endfirsthead
  \toprule
  Instance & ALNS Obj & ALNS Inf & VNS Obj & VNS Inf & $\Delta$ Obj \\
  \midrule
  \endhead
  \bottomrule
  \endfoot
  Center486\_G0 & 195 & 0{,}0000 & 201 & 0{,}0000 & $-$6 \\
  Center486\_G1 & 259 & 0{,}0000 & 268 & 0{,}0000 & $-$9 \\
  Center486\_G2 & 270 & 0{,}0000 & 283 & 0{,}0000 & $-$13 \\
  Center486\_G3 & 227 & 0{,}0000 & 222 & 0{,}0000 & $+$5 \\
  Center486\_G4 & 236 & 0{,}0000 & 175 & 0{,}0000 & $+$61 \\
  Center486\_G5 & 202 & 0{,}0000 & 172 & 0{,}0000 & $+$30 \\
  Center486\_G6 & 232 & 0{,}0000 & 212 & 0{,}0000 & $+$20 \\
  Center486\_G7 & 258 & 0{,}0000 & 210 & 0{,}0000 & $+$48 \\
  Center486\_G8 & 224 & 0{,}0000 & 194 & 0{,}0000 & $+$30 \\
  Center486\_G9 & 227 & 0{,}0000 & 251 & 0{,}0000 & $-$24 \\
  \midrule
  Corners486\_G0  & 277 & 0{,}0000 & 219 & 0{,}0000 & $+$58 \\
  Corners486\_G1  & 190 & 0{,}0000 & 183 & 0{,}0000 & $+$7 \\
  Corners486\_G2  & 203 & 0{,}0000 & 186 & 0{,}0000 & $+$17 \\
  Corners486\_G3  & 191 & 0{,}0000 & 167 & 0{,}0000 & $+$24 \\
  Corners486\_G4  & 273 & 0{,}0000 & 223 & 0{,}0119 & $+$50 \\
  Corners486\_G5  & 371 & 0{,}0000 & 252 & 0{,}0000 & $+$119 \\
  Corners486\_G6  & 315 & 0{,}0000 & 202 & 0{,}0000 & $+$113 \\
  Corners486\_G7  & 182 & 0{,}0099 & 174 & 0{,}0521 & $+$8 \\
  Corners486\_G8  & 244 & 0{,}0000 & 219 & 0{,}0101 & $+$25 \\
  Corners486\_G9  & 296 & 0{,}0000 & 192 & 0{,}0000 & $+$104 \\
  \midrule
  Diagonal486\_G0 & 347 & 0{,}0000 & 211 & 0{,}0197 & $+$136 \\
  Diagonal486\_G1 & 211 & 0{,}0000 & 232 & 0{,}0624 & $-$21 \\
  Diagonal486\_G2 & 230 & 0{,}0000 & 182 & 0{,}0736 & $+$48 \\
  Diagonal486\_G3 & 223 & 0{,}0000 & 195 & 0{,}0277 & $+$28 \\
  Diagonal486\_G4 & 270 & 0{,}0000 & 220 & 0{,}0077 & $+$50 \\
  Diagonal486\_G5 & 296 & 0{,}0000 & 190 & 0{,}0059 & $+$106 \\
  Diagonal486\_G6 & 294 & 0{,}0000 & 199 & 0{,}0633 & $+$95 \\
  Diagonal486\_G7 & 171 & 0{,}0000 & 196 & 0{,}0300 & $-$25 \\
  Diagonal486\_G8 & 174 & 0{,}0000 & 172 & 0{,}0000 & $+$2 \\
  Diagonal486\_G9 & 254 & 0{,}0000 & 199 & 0{,}0289 & $+$55 \\
\end{longtable}

\end{document}
